%%=============================================================================
%% Methodologie
%%=============================================================================

\chapter{\IfLanguageName{dutch}{Methodologie}{Methodology}}%
\label{ch:methodologie}

\paragraph{Fase 0: Onderzoeksvoorstel (12/12 - 23/01)}

Dit is de voorbereidende fase waar het onderwerp wordt afgebakend, de methodologie wordt uitgewerkt, een kleine literatuurstudie wordt uitgevoerd en de scope van het onderzoek wordt bepaald.

\paragraph{Fase 1: Literatuurstudie (23/01 - 15/02)}     
De literatuurstudie uit het onderzoeksvoorstel wordt waar nodig verder uitgebreid om de beste theorieën en oplossingen uit de bestaande wetenschappelijke lectuur te halen, die helpen de onderzoeksvraag te beantwoorden. Dit vormt de theoretische basis voor een mogelijk Proof of Concept en onderzoekt ook hoe vergelijkbare systemen omgaan met onboarding. 

\paragraph{Fase 2: Requirements verzamelen (15/02 - 01/03)}

In deze fase wordt gezocht naar de verschillende criteria en vereisten waaraan een goede oplossing moet voldoen. 
\\

Dit gebeurt in eerste instantie door overleg en interviews te voeren met verschillende stakeholders (zoals de copromotor en werknemers van 24flow) om het huidige onboardingsproces te definiëren en de gebreken, belangrijkste vereisten en flows te identificeren. Deze worden vervolgens geprioriteerd aan de hand van de MoSCoW-methode. 
\\

Uiteindelijk resulteert dit in een geordende lijst van requirements voor het prototype.
\\

Verder wordt in deze fase het platform verkend en de Salesforce documentatie doorgenomen om te bepalen welke native Salesforce-functionaliteiten het meest toepasbaar zijn voor de mogelijke oplossing. Daaruit volgt een shortlist van de geschikte tools.

\paragraph{Fase 3: Ontwerp prototype (01/03 - 15/03)}
Aan de hand van de opgestelde requirements en de Salesforce-analyse wordt een ontwerp opgesteld voor het onboarding prototype.
\\

Dit is een verzameling van verschillende use cases en wireframes die samen de blauwdruk vormen voor de mogelijke onboarding-oplossing. 
\\

Dit ontwerp wordt besproken met de promotor en vervolgens bijgewerkt naargelang de gekregen feedback, om zo tot een definitief ontwerp te komen voor de implementatie.

\paragraph{Fase 4: Implementatie prototype (15/03 - 15/05)}
Hier wordt het gemaakte ontwerp effectief geïmplementeerd binnen het 24flow-platform op een developer-omgeving. De belangrijkste requirements worden gebouwd met behulp van de geselecteerde Salesforce-functionaliteiten en Apex-logica.
\\

Deze fase leidt tot een werkend proof of concept dat een semi-geautomatiseerd onboarding- en guidance-systeem op het 24flow-platform voorziet.

\paragraph{Fase 5: Evaluatie (15/05 - 29/05)}
In deze laatste fase wordt het gerealiseerde proof of concept kritisch geëvalueerd om te bepalen in hoeverre het een antwoord biedt op de hoofdonderzoeksvraag.
\\

Allereerst wordt het prototype afgetoetst aan de vooropgestelde requirements, om te bepalen welke functionaliteiten succesvol zijn geïmplementeerd en waar er mogelijks is afgeweken van het originele ontwerp en waarom. 
\\

Vervolgens wordt de werking van de oplossing kwalitatief geëvalueerd. Dit gebeurt enerzijds door het prototype te spiegelen aan de theorie uit de literatuurstudie, en anderzijds door het resultaat te demonstreren aan interne 24flow-medewerkers. Via deze demonstraties wordt kwalitatieve feedback verzameld over criteria zoals de verwachte opzettijd en de mate waarin een omgeving met deze nieuwe tools zelfstandig geconfigureerd kan worden. Op basis van deze inzichten wordt ten slotte het uiteindelijke besluit, de algemene conclusie en de praktische meerwaarde van het onderzoek geformuleerd.